\problem{4.12}{Unnormalized hydrogen radial functions at $n=3$}
With $\rho=2r/(3\aB)$,
the hydrogen radial equation admits
$R_{n\ell}=C\,\rho^\ell\ee^{-\rho/2}
L_{n-\ell-1}^{2\ell+1}(\rho)$.
For $n=3$ the relevant associated Laguerre
polynomials are $L_2^1(\rho)=
\frac12(\rho^2-6\rho+6)$,
$L_1^3(\rho)=4-\rho$, and $L_0^5=1$.
Therefore, up to overall normalization,
\begin{align}
\boxed{R_{30}}&\propto
(6-6\rho+\rho^2)\ee^{-\rho/2},\\
\boxed{R_{31}}&\propto
\rho(4-\rho)\ee^{-\rho/2},\\
\boxed{R_{32}}&\propto
\rho^2\ee^{-\rho/2}.
\end{align}
The terminating polynomial degrees are respectively
$2,1,0$; all decay as $\ee^{-r/(3\aB)}$.

\problem{4.13}{Normalizing the $n=2$ hydrogen eigenfunctions}
For $\rho=r/\aB$, the radial functions take
the normalized forms
\begin{align}
\boxed{R_{20}(r)}&=\frac{1}{2\sqrt2\,\aB^{3/2}}
(2-\rho)\ee^{-\rho/2},\\
\boxed{R_{21}(r)}&=\frac1{2\sqrt6\,\aB^{3/2}}
\rho\ee^{-\rho/2}.
\end{align}
For example
$\int_0^\infty |R_{20}|^2r^2\dd r
=\frac18\int_0^\infty
\rho^2(2-\rho)^2\ee^{-\rho}\dd\rho
=\frac18(8-24+24)=1$.
For $R_{21}$, the integral equals
$\frac1{24}\int\rho^4\ee^{-\rho}\dd\rho=1$.
Combining with the normalized harmonics,
\begin{align}
\boxed{\psi_{200}}&=
\frac{(2-\rho)\ee^{-\rho/2}}{4\sqrt{2\pi}\,\aB^{3/2}},\\
\boxed{\psi_{210}}&=
\frac{\rho\cos\theta\,\ee^{-\rho/2}}
{4\sqrt{2\pi}\,\aB^{3/2}},\\
\boxed{\psi_{211}}&=
-\frac{\rho\sin\theta\,\ee^{\ii\phi}\ee^{-\rho/2}}
{8\sqrt{\pi}\,\aB^{3/2}},\\
\boxed{\psi_{21,-1}}&=
\frac{\rho\sin\theta\,\ee^{-\ii\phi}\ee^{-\rho/2}}
{8\sqrt{\pi}\,\aB^{3/2}}.
\end{align}
Angular orthonormality and the radial normalization
ensure unit three-dimensional norm.

\problem{4.14}{Laguerre polynomials and a higher radial solution}
With the ordinary Laguerre convention
$L_q(x)=\ee^x(q!)^{-1}
\dd^q(\ee^{-x}x^q)/\dd x^q$,
\begin{align}
L_0&=1,&L_1&=1-x,\\
L_2&=1-2x+\tfrac12x^2,&
L_3&=1-3x+\tfrac32x^2-\tfrac16x^3.
\end{align}
The radial polynomial for $n=5,\ell=2$ has
$n-\ell-1=2$ and parameter $2\ell+1=5$,
so
\begin{equation}
\boxed{L_2^5(\rho)=21-7\rho+\frac12\rho^2}.
\end{equation}
In the textbook's radial ansatz $u(\rho)=
\rho^{\ell+1}\ee^{-\rho/2}v(\rho)$,
the polynomial is
$\boxed{v(\rho)\propto21-7\rho+\rho^2/2}$.
Checking the series recurrence obtained by
substituting $v=\sum_jc_j\rho^j$ into
\begin{equation}
\rho v''+(2\ell+2-\rho)v'
+(n-\ell-1)v=0
\end{equation}
gives $(j+1)(j+2\ell+2)c_{j+1}=
(j+\ell+1-n)c_j$.
At $n=5,\ell=2$:
$c_1/c_0=-1/3$, $c_2/c_1=-1/14$,
$c_3=0$, confirming the polynomial.

\problem{4.15}{Hydrogen mean radii and Cartesian position moments}
The $1s$ density is
$|\psi_{100}|^2=\ee^{-2r/\aB}/(\pi\aB^3)$.
Using $\int_0^\infty r^j\ee^{-\beta r}\dd r
=j!/\beta^{j+1}$,
\begin{align}
\avg r&=\frac4{\aB^3}\int_0^\infty
r^3\ee^{-2r/\aB}\dd r
=\boxed{\tfrac32\aB},\\
\avg{r^2}&=\frac4{\aB^3}\int_0^\infty
r^4\ee^{-2r/\aB}\dd r
=\boxed{3\aB^2}.
\end{align}
Spherical symmetry implies
$\boxed{\avg x=0,\ \avg{x^2}=\aB^2}$.
For the $n=2,\ell=1,m=1$ state,
the standard radial integral gives
$\avg{r^2}=30\aB^2$.
Its angular density is proportional to
$\sin^2\theta$, for which
$\avg{\sin^2\theta\cos^2\phi}=2/5$;
hence
\begin{equation}
\boxed{\avg{x^2}_{211}=12\aB^2}.
\end{equation}

\problem{4.16}{Most probable hydrogenic ground-state radius}
The probability in a spherical shell of thickness
$\dd r$ is $P(r)\dd r$ with
\begin{equation}
P(r)=4\pi r^2|\psi_{100}(r)|^2
=\frac{4r^2}{\aB^3}\ee^{-2r/\aB}.
\end{equation}
Differentiate:
$P'(r)=\frac{8r}{\aB^3}\ee^{-2r/\aB}
(1-r/\aB)$.
The positive stationary point is
$r=\aB$; the density vanishes at $r=0$
and at infinity, so the maximum is
\begin{equation}
\boxed{r_{\rm most\ probable}=\aB}.
\end{equation}
The radial maximum is not the same
as the maximum of the spatial density,
which occurs at $r=0$.

\problem{4.17}{A double-commutator shortcut for $\langle zHz\rangle$}
Since $[H,z]=-\ii\sh p_z/m$ and
$[z,p_z]=\ii\sh$, one has
\begin{equation}
[z,[H,z]]=\frac{\sh^2}{m}I.
\end{equation}
Expanding the double commutator gives
$[z,[H,z]]=2zHz-z^2H-Hz^2$.
In the normalized ground eigenstate
$H\ket{100}=E_1\ket{100}$,
\begin{equation}
2\avg{zHz}-2E_1\avg{z^2}
=\frac{\sh^2}{m}.
\end{equation}
By Problem 4.15, $\avg{z^2}=\aB^2$ and
$E_1=-\sh^2/(2m\aB^2)$
(fixed-nucleus approximation), so
\begin{equation}
\boxed{\avg{zHz}=E_1\aB^2+\frac{\sh^2}{2m}=0}.
\end{equation}
No six-dimensional integral is necessary.

\problem{4.18}{Stationarity within the hydrogen $n=2$ eigenspace}
For the initial state
$\Psi(\mathbf r,0)=
(\psi_{211}+\psi_{21,-1})/\sqrt2$,
both components share energy
$E_2=-13.6/4\,{\rm eV}=-3.40\,{\rm eV}$.
Therefore
\begin{equation}
\boxed{\Psi(\mathbf r,t)=
\ee^{-\ii E_2t/\sh}\,
\frac{\psi_{211}(\mathbf r)+\psi_{21,-1}(\mathbf r)}
{\sqrt2}}.
\end{equation}
Using $Y_{11}\propto-\sin\theta\ee^{\ii\phi}$,
$Y_{1,-1}\propto+\sin\theta\ee^{-\ii\phi}$,
the angular part is proportional to
$-\ii\sin\theta\sin\phi$ (a $p_y$ orbital).
Only a global time-dependent phase appears:
$|\Psi(\mathbf r,t)|^2$ is stationary.
For Coulomb $V\propto1/r$, the virial theorem
gives $\avg V=2E_2$; thus
\begin{equation}
\boxed{\avg V=2E_2=-6.80\,{\rm eV}},
\end{equation}
independent of time.

\problem{4.19}{Hydrogenic scaling with nuclear charge $Z$}
For the fixed-nucleus Coulomb potential
$V_Z=-Ze^2/(4\pi\epsilon_0r)$,
replace $e^2$ by $Ze^2$ in the hydrogen
radius and energy formulas. Hence
\begin{equation}
\boxed{a_Z=\frac{\aB}{Z},\quad
E_n(Z)=-\frac{13.6\,Z^2}{n^2}\,{\rm eV},\quad
|E_1(Z)|=13.6Z^2\,{\rm eV},\quad
R_Z=Z^2R_H}.
\end{equation}
With finite nuclear mass, additionally
replace electron mass by the corresponding
electron--nucleus reduced mass.
Lyman wavelengths scale as $1/Z^2$:
for $Z=2$, the series is at one quarter
of the hydrogen wavelengths (extreme UV);
for $Z=3$, one ninth (extreme UV/soft-X-ray
boundary for the shortest wavelengths).
The classification depends on the exact
wavelength threshold adopted.

\problem{4.20}{Gravitational hydrogen analogy for the Earth and Sun}
For Earth mass $m_E$, solar mass $M$, and
reduced mass $\mu=m_EM/(m_E+M)$,
\begin{equation}
\boxed{V(r)=-\frac{GMm_E}{r},\quad
a_g=\frac{\sh^2}{\mu GMm_E},\quad
E_n=-\frac{\mu(GMm_E)^2}{2\sh^2n^2}}.
\end{equation}
Using $M\simeq1.9885\times10^{30}$ kg,
$m_E\simeq5.9722\times10^{24}$ kg,
$\sh=1.05457\times10^{-34}$ J s, and
$G\simeq6.6743\times10^{-11}$ SI units,
\begin{equation}
\boxed{a_g\simeq2.35\times10^{-138}\ {\rm m}}.
\end{equation}
A classical circular orbit at radius $r_0$ has
$E_{\rm cl}=-GMm_E/(2r_0)$; comparing with $E_n$
gives $r_0=a_gn^2$ and
\begin{equation}
\boxed{n=\sqrt{r_0/a_g}\simeq
2.52\times10^{74}}
\end{equation}
for one astronomical unit.
For $n\gg1$, the next-level energy loss is
$\Delta E\sim2|E_n|/n$, giving
$\boxed{\Delta E\simeq2.10\times10^{-41}\ {\rm J}}$.
The corresponding quantum wavelength
$\lambda=hc/\Delta E$ is
\begin{equation}
\boxed{\lambda\simeq9.46\times10^{15}\ {\rm m}
\simeq1.00\ {\rm light\ year}}.
\end{equation}
The near-one-light-year number reflects
the numerical relationship of the
Earth--Sun orbital parameters and
fundamental constants; the toy model
is not a realistic description of
gravitational radiation from the orbit.
